\documentclass[a4paper,11pt]{article}


\usepackage[margin=2.2cm]{geometry}
\usepackage{graphicx}
\usepackage{float}
\usepackage{booktabs}
\usepackage{longtable}
\usepackage{array}
\usepackage{xcolor}
\usepackage{hyperref}
\usepackage{enumitem}
\usepackage{amsmath}

\hypersetup{
	colorlinks=true,
	linkcolor=black,
	urlcolor=blue
}

\setcounter{secnumdepth}{0}

\title{Missing Values and Outlier Detection}
\author{Giuseppe Rudi}
\date{\today}

\begin{document}
	
	\maketitle
	
	
	\section{Introduction}

	
	This document describes the focused analysis of missing values and outliers applied to the Formula 1 reconciled database.
	
	This phase is separated from the General Data Quality Assessment. Some data quality problems require a more specific interpretation.
	This is especially true for missing values and numerical outliers.
	In Formula 1 data, a missing value is not always an error.
	Similarly, an extreme lap time or speed value is not always wrong.
	It may be caused by pit stops, Safety Car periods, deleted laps, inaccurate laps, rain, crashes, red flags, or other racing situations.
	
	
	For this reason, this focused analysis is designed as a separate step after the General DQA and before the final Data Cleaning phase.
	
	The goal is not to clean the data automatically.
	The goal is to classify missing values, flag suspicious numerical values, and produce evidence for later cleaning or analytical filtering decisions.
	
	% ============================================================
	\subsection{Position in the Project Pipeline}
	% ============================================================
	
	The focused missing values and outlier detection phase is placed after the General DQA.
	
	\begin{center}
		\begin{tabular}{c}
			Raw and external sources \\
			$\downarrow$ \\
			Extraction and re-engineering \\
			$\downarrow$ \\
			PostgreSQL reconciled database \\
			$\downarrow$ \\
			General Data Quality Assessment \\
			$\downarrow$ \\
			\textbf{Missing Values and Outlier Detection} \\
			$\downarrow$ \\
			Data Cleaning \\
			$\downarrow$ \\
			ETL toward the Data Warehouse \\
			$\downarrow$ \\
			Star schema and OLAP analysis
		\end{tabular}
	\end{center}
	

	% ============================================================
	\subsection{Scope of the Focused Analysis}
	% ============================================================
	
	The focused analysis is applied only to the two most important and most complex tables:
	
	\begin{itemize}[leftmargin=1.2cm]
		\item \texttt{lap}
		\item \texttt{result}
	\end{itemize}
	
	These tables are selected because they contain the main analytical data used in the project.



	% ============================================================
	\section{Missing Values}
	% ============================================================
	
	The goal of the focused missing values analysis is not only to count null values.
	The goal is to interpret why a value is missing and how much this missing value affects the analytical use of the row.
	

	
	\subsubsection{Missing Value Classes}
	
	The script uses two main missing value classes.
	
	\begin{longtable}{p{4cm}p{6cm}p{5.5cm}}
		\toprule
		\textbf{Class} & \textbf{Meaning} & \textbf{Example} \\
		\midrule
		\endfirsthead
		
		\toprule
		\textbf{Class} & \textbf{Meaning} & \textbf{Example} \\
		\midrule
		\endhead
		
		\texttt{EXPLAINED\_NULL}
		& The value is missing, but the missingness can be explained by a clear Formula 1 domain rule or by another observed attribute in the row.
		& \texttt{q2\_ms} is missing because the driver was eliminated in Q1; \texttt{pit\_in\_time\_ms} is missing because the lap is not an in-lap. \\
		\midrule
		
		\texttt{SUSPICIOUS\_NULL}
		& The value is missing and the available domain rules do not provide a clear explanation.
		& \texttt{lap\_time\_ms} is missing even though the lap is not a pit-in lap, not a pit-out lap, not deleted, and marked as accurate. \\
		\bottomrule
	\end{longtable}
		
	
	\subsubsection{Missing Information Area}
	
	In addition to the missing value class, the script assigns a second field called
	\texttt{missing\_information\_area}.
	
	The goal of this field is to describe which type of analytical information is missing from the row.
	
	This choice is more useful than a generic impact level because the same missing value can affect different analyses in different ways.
	
	For example, a lap with a missing speed value may still be usable for lap time analysis, but it cannot be used for speed-based analysis.
	
	Similarly, a lap with a missing compound may still contain a lap time, but it cannot be used for tyre-based or fair contextual comparison.
	
	Therefore, the script does not directly decide which dashboard or analysis must exclude the row.
	
	Instead, it records the type of missing information.
	The later visualization phase can decide how to use or filter the
	row depending on the analytical goal.
	
	\begin{longtable}{p{5cm}p{9.5cm}}
		\toprule
		\textbf{Missing Information Area} & \textbf{Meaning} \\
		\midrule
		\endfirsthead
		
		\toprule
		\textbf{Missing Information Area} & \textbf{Meaning} \\
		\midrule
		\endhead
		
		\texttt{NONE}
		& The missing value does not remove useful analytical information. This is used for expected nulls. \\
		\midrule
		
		\texttt{LAP\_TIME\_INFORMATION}
		& The missing value affects the main lap time information. This is used when \texttt{lap\_time\_ms} is missing. \\
		\midrule
		
		\texttt{SECTOR\_INFORMATION}
		& The missing value affects sector-based performance information. This is used when one or more sector times are missing. \\
		\midrule
		
		\texttt{SPEED\_INFORMATION}
		& The missing value affects speed-based information. This is used when one or more speed attributes are missing. \\
		\midrule
		
		\texttt{TYRE\_INFORMATION}
		& The missing value affects tyre-related information. This is used when attributes such as \texttt{compound} or \texttt{tyre\_life} are missing. \\
		\midrule
		
		\texttt{WEATHER\_INFORMATION}
		& The missing value affects weather context information. This can be used when a lap cannot be associated with reliable weather information. \\
		\midrule
		
		\texttt{TRACK\_STATUS\_INFORMATION}
		& The missing value affects track condition information. This can be used when the track status context is missing or not interpretable. \\
		\midrule
		
		\texttt{PIT\_INFORMATION}
		& The missing value affects pit-related information. This can be used when pit information is required to interpret special laps. \\

		\bottomrule
	\end{longtable}
	
	For the \texttt{result} table, the following information areas are used:
	
	\begin{longtable}{p{5cm}p{9.5cm}}
		\toprule
		\textbf{Missing Information Area} & \textbf{Meaning} \\
		\midrule
		\endfirsthead
		
		\toprule
		\textbf{Missing Information Area} & \textbf{Meaning} \\
		\midrule
		\endhead
		
		\texttt{NONE}
		& The missing value does not remove useful analytical information in the current session context.
		This is used for expected nulls, such as qualifying times in Race sessions or race-related attributes in Qualifying sessions. \\
		\midrule
		
		\texttt{CLASSIFICATION\_INFORMATION}
		& The missing value affects the classification information of the driver in the session.
		This is used when \texttt{position} is missing, or when \texttt{classified\_position} is missing in a Race session.
		This area is considered critical because classification information is necessary to interpret the final order or result of the driver. \\
		\midrule
		
		\texttt{QUALIFYING\_INFORMATION}
		& The missing value affects qualifying performance information.
		This is used when \texttt{q1\_ms}, \texttt{q2\_ms}, or \texttt{q3\_ms} are missing in a context where they should normally be present according to the qualifying progression. \\
		\midrule
		
		\texttt{RACE\_CONTEXT\_INFORMATION}
		& The missing value affects additional race context information.
		This is used for missing race-related attributes such as \texttt{grid\_position}, \texttt{status}, \texttt{points}, \texttt{laps}, or \texttt{time\_ms}.
		These attributes may be useful for specific analyses, but they are not always essential for the main lap performance analysis of the project. \\
		\bottomrule
	\end{longtable}
	
	The final interpretation of a missing value is obtained by combining the two fields:
	
	\begin{center}
		\begin{tabular}{lll}
			\toprule
			\textbf{Column} & \textbf{Missing Class} & \textbf{Missing Information Area} \\
			\midrule
			\texttt{q3\_ms} & \texttt{EXPLAINED\_NULL} & \texttt{NONE} \\
			\texttt{sector1\_time\_ms} & \texttt{EXPLAINED\_NULL} & \texttt{SECTOR\_INFORMATION} \\
			\texttt{compound} & \texttt{SUSPICIOUS\_NULL} & \texttt{TYRE\_INFORMATION} \\
			\texttt{lap\_time\_ms} & \texttt{SUSPICIOUS\_NULL} & \texttt{LAP\_TIME\_INFORMATION} \\
			\bottomrule
		\end{tabular}
	\end{center}
	
	This approach keeps the classification simple but more useful for the later analytical
	steps.
	
	The field \texttt{missing\_class} explains why the value is null.
	The field \texttt{missing\_information\_area} explains what kind of information is
	missing from the row.
	
	For example, if a later visualization requires tyre information, it can exclude rows where \texttt{missing\_information\_area} contains \texttt{TYRE\_INFORMATION}.
	If a visualization only compares basic lap times, it may still use rows where speed
	or sector information is missing, as long as the lap time itself is present.

	
	
	% ============================================================
	\section{Missing Values on the Result Table}
	% ============================================================
	
	The \texttt{result} table contains one row for one driver in one session.
	It includes both qualifying-related attributes and race-related attributes.
	
	\begin{table}[H]
		\centering
		\small
		\begin{tabular}{p{4cm} p{10.5cm}}
			\toprule
			\textbf{Category} & \textbf{Attributes} \\
			\midrule
			Qualifying Attributes &
			\texttt{q1\_ms}, \texttt{q2\_ms}, \texttt{q3\_ms} \\[0.2cm]
			
			Race-related Attributes &
			\texttt{classified\_position}, \texttt{status}, \texttt{points}, \texttt{laps}, \texttt{time\_ms} \\
			\bottomrule
		\end{tabular}
	\end{table}
	
	Because the meaning of these attributes depends on the session type, missing values in this table are analyzed differently for Qualifying and Race sessions.
	
	The script uses the \texttt{session\_type} associated with each \texttt{session\_id}
	to decide which rules must be applied.
	
	% ------------------------------------------------------------
	\subsection{Race Sessions}
	% ------------------------------------------------------------
	
	For Race sessions, qualifying time attributes are not applicable.
	
	Therefore, missing values in \texttt{q1\_ms}, \texttt{q2\_ms}, and \texttt{q3\_ms}
	are classified as explained nulls with no missing analytical information.
	
	\begin{center}
		\begin{tabular}{lll}
			\toprule
			\textbf{Missing Attribute} & \textbf{Missing Class} & \textbf{Missing Information Area} \\
			\midrule
			\texttt{q1\_ms} & \texttt{EXPLAINED\_NULL} & \texttt{NONE} \\
			\texttt{q2\_ms} & \texttt{EXPLAINED\_NULL} & \texttt{NONE} \\
			\texttt{q3\_ms} & \texttt{EXPLAINED\_NULL} & \texttt{NONE} \\
			\bottomrule
		\end{tabular}
	\end{center}
	
	These values are explained by the session type: a Race session does not contain Q1,
	Q2, or Q3 segment times. Therefore, no useful race-result information is missing.
	
	By contrast, race-related attributes such as \texttt{classified\_position},
	\texttt{status}, \texttt{points}, \texttt{laps}, and \texttt{time\_ms} are expected
	to be present in Race sessions. If one of these values is missing, the missing value
	is classified as suspicious because race-result information is incomplete.
	
	\begin{longtable}{p{5cm}p{4cm}p{6cm}}
		\toprule
		\textbf{Missing Attribute} & \textbf{Missing Class} & \textbf{Missing Information Area} \\
		\midrule
		\endfirsthead
		
		\toprule
		\textbf{Missing Attribute} & \textbf{Missing Class} & \textbf{Missing Information Area} \\
		\midrule
		\endhead
		
		\texttt{position}
		& \texttt{SUSPICIOUS\_NULL}
		& \texttt{CLASSIFICATION\_INFORMATION} \\
		\midrule
		
		\texttt{classified\_position}
		& \texttt{SUSPICIOUS\_NULL}
		& \texttt{CLASSIFICATION\_INFORMATION} \\
		\midrule
		
		\texttt{grid\_position}
		& \texttt{SUSPICIOUS\_NULL}
		& \texttt{RACE\_CONTEXT\_INFORMATION} \\
		\midrule
		
		\texttt{status}
		& \texttt{SUSPICIOUS\_NULL}
		& \texttt{RACE\_CONTEXT\_INFORMATION} \\
		\midrule
		
		\texttt{points}
		& \texttt{SUSPICIOUS\_NULL}
		& \texttt{RACE\_CONTEXT\_INFORMATION} \\
		\midrule
		
		\texttt{laps}
		& \texttt{SUSPICIOUS\_NULL}
		& \texttt{RACE\_CONTEXT\_INFORMATION} \\
		\midrule
		
		\texttt{time\_ms}
		& \texttt{SUSPICIOUS\_NULL}
		& \texttt{RACE\_CONTEXT\_INFORMATION} \\
		\bottomrule
	\end{longtable}

	
	% ------------------------------------------------------------
	\subsection{Qualifying Sessions}
	% ------------------------------------------------------------
	

	In Qualifying sessions, \texttt{position} is fundamental because it is used to interpret the qualifying progression.
	Without \texttt{position}, the script cannot reliably decide whether a missing \texttt{q2\_ms} or \texttt{q3\_ms} is expected or suspicious.
	
	\begin{longtable}{p{4cm}p{4cm}p{6cm}}
		\toprule
		\textbf{Missing Attribute} & \textbf{Missing Class} & \textbf{Missing Information Area} \\
		\midrule
		\endfirsthead
		
		\toprule
		\textbf{Missing Attribute} & \textbf{Missing Class} & \textbf{Missing Information Area} \\
		\midrule
		\endhead
		
		\texttt{position}
		& \texttt{SUSPICIOUS\_NULL}
		& \texttt{CLASSIFICATION\_INFORMATION} \\
		\bottomrule
	\end{longtable}
	
	For Qualifying sessions, the attributes \texttt{q1\_ms}, \texttt{q2\_ms}, and \texttt{q3\_ms} must be interpreted according to the qualifying progression.
	
	Not all drivers are expected to have Q2 or Q3 times.
	Drivers eliminated in Q1 should normally have no Q2 or Q3 time.
	Drivers eliminated in Q2 should normally have no Q3 time.
	
	For this reason, missing qualifying times can be either explained or suspicious,
	depending on the final qualifying position.
	
	\begin{longtable}{p{4.5cm}p{2.2cm}p{3.5cm}p{4cm}}
		\toprule
		\textbf{Condition} & \textbf{Missing Attribute} & \textbf{Missing Class} & \textbf{Missing Information Area} \\
		\midrule
		\endfirsthead
		
		\toprule
		\textbf{Condition} & \textbf{Missing Attribute} & \textbf{Missing Class} & \textbf{Missing Information Area} \\
		\midrule
		\endhead
		
		\texttt{position > 15}
		& \texttt{q2\_ms}
		& \texttt{EXPLAINED\_NULL}
		& \texttt{NONE} \\
		\midrule
		
		\texttt{position > 15}
		& \texttt{q3\_ms}
		& \texttt{EXPLAINED\_NULL}
		& \texttt{NONE} \\
		\midrule
		
		\texttt{11 <= position <= 15}
		& \texttt{q3\_ms}
		& \texttt{EXPLAINED\_NULL}
		& \texttt{NONE} \\
		\bottomrule
	\end{longtable}
	
	These missing values are not considered data quality errors because they are explained by the structure of the qualifying session. In these cases, no useful qualifying information is missing, because the driver did not reach the corresponding qualifying phase.
	
	However, if a qualifying time is missing even though the driver's position indicates that the driver should normally have reached that phase, the missing value is classified
	as suspicious and the missing information area is \texttt{QUALIFYING\_INFORMATION}.
	
	\begin{longtable}{p{4cm}p{2.2cm}p{3.3cm}p{4.4cm}}
		\toprule
		\textbf{Condition} & \textbf{Missing Attribute} & \textbf{Missing Class} & \textbf{Missing Information Area} \\
		\midrule
		\endfirsthead
		
		\toprule
		\textbf{Condition} & \textbf{Missing Attribute} & \textbf{Missing Class} & \textbf{Missing Information Area} \\
		\midrule
		\endhead
		
		Qualifying session
		& \texttt{q1\_ms}
		& \texttt{SUSPICIOUS\_NULL}
		& \texttt{QUALIFYING\_INFORMATION} \\
		\midrule
		
		\texttt{position <= 15}
		& \texttt{q2\_ms}
		& \texttt{SUSPICIOUS\_NULL}
		& \texttt{QUALIFYING\_INFORMATION} \\
		\midrule
		
		\texttt{position <= 10}
		& \texttt{q3\_ms}
		& \texttt{SUSPICIOUS\_NULL}
		& \texttt{QUALIFYING\_INFORMATION} \\
		\bottomrule
	\end{longtable}
	
	In these cases, qualifying information is incomplete for that row.
	For example, if a driver is classified in the top 10 but \texttt{q3\_ms} is missing,
	the script flags the value as \texttt{SUSPICIOUS\_NULL} with
	\texttt{QUALIFYING\_INFORMATION}.
	
	The row is not corrected automatically, because special cases may exist, such as
	drivers who did not set a valid time, deleted laps, penalties, or session anomalies.
	The script only records that qualifying information is missing and that the case
	requires review.
	
	\vspace{2em}
	
	For Qualifying sessions, race-specific attributes such as \texttt{points},
	\texttt{laps}, \texttt{classified\_position}, \texttt{status}, and \texttt{time\_ms}
	are not central to the qualifying performance analysis.
	
	If these attributes are missing in a Qualifying session, the missing values are treated
	as explained by the session type and no useful qualifying information is considered
	missing.
	
	\begin{center}
		\begin{tabular}{lll}
			\toprule
			\textbf{Missing Attribute} & \textbf{Missing Class} & \textbf{Missing Information Area} \\
			\midrule
			\texttt{classified\_position} & \texttt{EXPLAINED\_NULL} & \texttt{NONE} \\
			\texttt{status} & \texttt{EXPLAINED\_NULL} & \texttt{NONE} \\
			\texttt{points} & \texttt{EXPLAINED\_NULL} & \texttt{NONE} \\
			\texttt{laps} & \texttt{EXPLAINED\_NULL} & \texttt{NONE}\\
			\texttt{time\_ms} & \texttt{EXPLAINED\_NULL} & \texttt{NONE} \\
			\texttt{grid\_position} & \texttt{EXPLAINED\_NULL} & \texttt{NONE} \\
			\bottomrule
		\end{tabular}
	\end{center}
	
	This rule keeps the analysis simple and avoids flagging expected null values as
	quality problems.
	

	% ============================================================
	\section{Missing Values on the Lap Table}
	% ============================================================
	
	The \texttt{lap} table is the most complex table of the reconciled database.
	Each row represents one lap performed by one driver in one session.
	
	This table contains several analytical attributes that are important for lap performance analysis:

	\begin{table}[H]
		\centering
		\small
		\begin{tabular}{p{4.2cm}p{4.2cm}p{4.2cm}}
			\toprule
			\textbf{Timing attributes} & \textbf{Speed attributes} & \textbf{Tyre / pit attributes} \\
			\midrule
			\texttt{lap\_time\_ms}     & \texttt{speed\_i1} & \texttt{compound} \\
			\texttt{sector1\_time\_ms} & \texttt{speed\_i2} & \texttt{tyre\_life} \\
			\texttt{sector2\_time\_ms} & \texttt{speed\_fl} & \texttt{pit\_in\_time\_ms} \\
			\texttt{sector3\_time\_ms} & \texttt{speed\_st} & \texttt{pit\_out\_time\_ms} \\
			\bottomrule
		\end{tabular}
	\end{table}
	
	% ------------------------------------------------------------
	\subsection{Pit-Related Missing Values}
	% ------------------------------------------------------------
	
	The attributes \texttt{pit\_in\_time\_ms} and \texttt{pit\_out\_time\_ms} represent session timestamps.
	
	If \texttt{pit\_in\_time\_ms} is present, the lap is an in-lap.
	If \texttt{pit\_out\_time\_ms} is present, the lap is an out-lap.
	
	These attributes are normally null in clean laps where the driver does not enter or exit the pit lane.
	Therefore, missing values in these two columns are usually expected and do not represent data quality problems.
	
	\begin{longtable}{p{4cm}p{4.5cm}p{3.5cm}p{4cm}}
		\toprule
		\textbf{Missing Attribute} & \textbf{Condition} & \textbf{Missing Class} & \textbf{Missing Information Area} \\
		\midrule
		\endfirsthead
		
		\toprule
		\textbf{Missing Attribute} & \textbf{Condition} & \textbf{Missing Class} & \textbf{Missing Information Area} \\
		\midrule
		\endhead
		
		\texttt{pit\_in\_time\_ms}
		& The lap is not an in-lap.
		& \texttt{EXPLAINED\_NULL}
		& \texttt{NONE} \\
		\midrule
		
		\texttt{pit\_out\_time\_ms}
		& The lap is not an out-lap.
		& \texttt{EXPLAINED\_NULL}
		& \texttt{NONE} \\
		\midrule
		
		\texttt{deleted\_reason}
		& \texttt{deleted = false}
		& \texttt{EXPLAINED\_NULL}
		& \texttt{NONE} \\
		\bottomrule
	\end{longtable}
	
	% ------------------------------------------------------------
	\subsection{Core Lap Performance Missing Values}
	% ------------------------------------------------------------
	
	The attribute \texttt{lap\_time\_ms} is the core measure of lap performance.
	If this value is missing, the row lacks the main lap time information.

	
	\begin{longtable}{p{4.5cm}p{4cm}p{3.5cm}p{4cm}}
		\toprule
		\textbf{Condition} & \textbf{Missing Attribute} & \textbf{Missing Class} & \textbf{Missing Information Area} \\
		\midrule
		\endfirsthead
		
		\toprule
		\textbf{Condition} & \textbf{Missing Attribute} & \textbf{Missing Class} & \textbf{Missing Information Area} \\
		\midrule
		\endhead
		
		
		No deletion or inaccuracy explanation is available
		& \texttt{lap\_time\_ms}
		& \texttt{SUSPICIOUS\_NULL}
		& \texttt{LAP\_TIME\_INFORMATION} \\
		\bottomrule
	\end{longtable}
	

	% ------------------------------------------------------------
	\subsection{Sector Time Missing Values}
	% ------------------------------------------------------------
	
	Sector times are important for sector-based performance analysis.
	However, missing sector times may occur in special laps.
	
	The script applies a conservative interpretation.
	A missing sector time is classified as \texttt{EXPLAINED\_NULL} only when there is a direct domain explanation.
	
	In this project, the only direct pit-related explanation used for sector times is the presence of \texttt{pit\_in\_time\_ms} for a missing \texttt{sector3\_time\_ms}.
	This is because \texttt{pit\_in\_time\_ms} means that the driver entered the pit lane during that lap, and the final part of the lap may be affected by pit entry or by the way the source records in-laps.
	
	By contrast, \texttt{pit\_out\_time\_ms} is not used as an automatic explanation for a missing \texttt{sector1\_time\_ms}.
	Even if the driver exits the pit lane during the lap, the first sector should not be considered missing by default only because of the pit exit information.
	
	
	\begin{longtable}{p{4.5cm}p{4cm}p{3.5cm}p{4cm}}
		\toprule
		\textbf{Condition} & \textbf{Missing Attribute} & \textbf{Missing Class} & \textbf{Missing Information Area} \\
		\midrule
		\endfirsthead
		
		\toprule
		\textbf{Condition} & \textbf{Missing Attribute} & \textbf{Missing Class} & \textbf{Missing Information Area} \\
		\midrule
		\endhead
		
		\texttt{pit\_in\_time\_ms} is present
		& \texttt{sector3\_time\_ms}
		& \texttt{EXPLAINED\_NULL}
		& \texttt{SECTOR\_INFORMATION} \\
		\midrule
		
		No direct sector-specific explanation is available
		& Any missing sector time
		& \texttt{SUSPICIOUS\_NULL}
		& \texttt{SECTOR\_INFORMATION} \\
		\bottomrule
	\end{longtable}
	
	A missing sector time always means that sector information is incomplete.
	The difference between \texttt{EXPLAINED\_NULL} and \texttt{SUSPICIOUS\_NULL} only depends on whether the script can find a direct and specific explanation.
	
	% ------------------------------------------------------------
	\subsection{Speed Missing Values}
	% ------------------------------------------------------------
	
	Speed attributes such as \texttt{speed\_i1}, \texttt{speed\_i2}, \texttt{speed\_fl}, and \texttt{speed\_st} are useful for speed-based analysis.
	
	A missing speed value does not necessarily make the whole lap unusable, but it means that speed information is incomplete.
	
	The script applies a conservative domain interpretation.
	Pit-related information is not used to explain all speed missing values.
	It is used only for the speed attributes that can reasonably be affected by pit entry, pit exit, or the start/end area of the lap.
	
	In this project, when \texttt{pit\_in\_time\_ms} or \texttt{pit\_out\_time\_ms} is present, missing values in \texttt{speed\_i1}, \texttt{speed\_fl}, and \texttt{speed\_st} can be considered explained.
	By contrast, missing \texttt{speed\_i2} is not automatically explained by pit information, because it is not directly associated with the pit entry or pit exit interpretation used in this simplified rule set.
	
	\begin{longtable}{p{4cm}p{5cm}p{3cm}p{4cm}}
		\toprule
		\textbf{Condition} & \textbf{Missing Attribute} & \textbf{Missing Class} & \textbf{Missing Information Area} \\
		\midrule
		\endfirsthead
		
		\toprule
		\textbf{Condition} & \textbf{Missing Attribute} & \textbf{Missing Class} & \textbf{Missing Information Area} \\
		\midrule
		\endhead
		
		\texttt{pit\_in\_time\_ms} is present
		& \texttt{speed\_i1}, \texttt{speed\_fl}, or \texttt{speed\_st}
		& \texttt{EXPLAINED\_NULL}
		& \texttt{SPEED\_INFORMATION} \\
		\midrule
		
		\texttt{pit\_out\_time\_ms} is present
		& \texttt{speed\_i1}, \texttt{speed\_fl}, or \texttt{speed\_st}
		& \texttt{EXPLAINED\_NULL}
		& \texttt{SPEED\_INFORMATION} \\
		\midrule
		
		
		Only pit information is present
		& \texttt{speed\_i2}
		& \texttt{SUSPICIOUS\_NULL}
		& \texttt{SPEED\_INFORMATION} \\
		\midrule
		
		No pit-related explanation, not deleted, and accurate lap
		& Any missing speed attribute
		& \texttt{SUSPICIOUS\_NULL}
		& \texttt{SPEED\_INFORMATION} \\
		\bottomrule
	\end{longtable}
	
	This rule avoids over-explaining speed missing values.
	The presence of an in-lap or out-lap is used as a possible explanation only for selected speed attributes.
	Missing \texttt{speed\_i2} is treated more cautiously and is classified as suspicious unless another explanation, such as \texttt{deleted = true} or \texttt{is\_accurate = false}, is available.

	
	% ------------------------------------------------------------
	\subsection{Tyre and Compound Missing Values}
	% ------------------------------------------------------------
	
	The attribute \texttt{compound} is essential in this project because lap performance cannot be interpreted fairly without knowing the tyre compound.
	
	Two laps with the same \texttt{lap\_time\_ms} may have very different meanings if they were set on different compounds.
	For this reason, a missing \texttt{compound} means that tyre information is incomplete.
	
	\begin{longtable}{p{4.5cm}p{4cm}p{3.5cm}p{4cm}}
		\toprule
		\textbf{Condition} & \textbf{Missing Attribute} & \textbf{Missing Class} & \textbf{Missing Information Area} \\
		\midrule
		\endfirsthead
		
		\toprule
		\textbf{Condition} & \textbf{Missing Attribute} & \textbf{Missing Class} & \textbf{Missing Information Area} \\
		\midrule
		\endhead
		
		Any lap
		& \texttt{compound}
		& \texttt{SUSPICIOUS\_NULL}
		& \texttt{TYRE\_INFORMATION} \\
		\midrule
		
		Any lap
		& \texttt{tyre\_life}
		& \texttt{SUSPICIOUS\_NULL}
		& \texttt{TYRE\_INFORMATION} \\
		\bottomrule
	\end{longtable}
	
	The value \texttt{tyre\_life} is important for tyre degradation and race pace comparison.
	If \texttt{tyre\_life} is missing, the row may still contain lap time and compound information, but tyre degradation information is incomplete.
	
	% ------------------------------------------------------------
	\subsection{Track Status and Weather Context Missing Values}
	% ------------------------------------------------------------
	
	Track status and weather information are important contextual elements for fair lap comparison.
	
	In the current reconciled \texttt{lap} table, \texttt{track\_status} is already available as a lap-level attribute.
	If this value is missing or not interpretable, the row lacks track condition information.
	
	Weather information is stored in a separate \texttt{weather} table.
	In this focused missing value script, weather alignment is not fully evaluated at lap level.
	However, if future ETL steps derive lap-level weather attributes and those values are missing, the corresponding missing information area will be \texttt{WEATHER\_INFORMATION}.
	
	\begin{longtable}{p{4.5cm}p{4cm}p{3.5cm}p{4cm}}
		\toprule
		\textbf{Condition} & \textbf{Missing Attribute} & \textbf{Missing Class} & \textbf{Missing Information Area} \\
		\midrule
		\endfirsthead
		
		\toprule
		\textbf{Condition} & \textbf{Missing Attribute} & \textbf{Missing Class} & \textbf{Missing Information Area} \\
		\midrule
		\endhead
		
		Any lap
		& \texttt{track\_status}
		& \texttt{SUSPICIOUS\_NULL}
		& \texttt{TRACK\_STATUS\_INFORMATION} \\
		\midrule
		
		Future derived weather context is missing
		& Weather-derived lap attribute
		& \texttt{SUSPICIOUS\_NULL}
		& \texttt{WEATHER\_INFORMATION} \\
		\bottomrule
	\end{longtable}
	
	This keeps the current script simple while leaving space for future weather-derived attributes.
	

	% ============================================================
	\section{Outlier Detection}
	% ============================================================
	
	Outlier detection is used to identify numerical values that are very different from the majority of values.
	
	In this project, outliers are not automatically considered errors.
	Formula 1 data naturally contains extreme values caused by race context.
	
	For example:
	
	\begin{itemize}[leftmargin=1.2cm]
		\item a slow sector may be caused by yellow flags;
		\item a low speed value may be caused by traffic or Safety Car;
		\item a very high lap time may be caused by rain, red flags, or an incomplete lap.
	\end{itemize}
	
	Therefore, outlier detection is used only to flag suspicious values for review.
	
	% ------------------------------------------------------------
	\subsection{Methods Used}
	% ------------------------------------------------------------
	
	The script uses two simple and interpretable univariate methods:
	
	\begin{itemize}[leftmargin=1.2cm]
		\item IQR method
		\item Modified Z-score
	\end{itemize}
	
	
	% ------------------------------------------------------------
	\subsubsection{IQR Method}
	% ------------------------------------------------------------
	
	The Interquartile Range method is based on the first quartile, the third quartile, and the interquartile range.
	
	\[
	IQR = Q3 - Q1
	\]
	
	The lower and upper bounds are defined as:
	
	\[
	LowerBound = Q1 - k \cdot IQR
	\]
	
	\[
	UpperBound = Q3 + k \cdot IQR
	\]
	
	A value outside this interval is flagged as an outlier.
	
	In this project, the IQR method is useful because it does not require a normal distribution assumption.
	
	% ------------------------------------------------------------
	\subsubsection{Modified Z-score}
	% ------------------------------------------------------------
	
	The Modified Z-score is based on the median and the Median Absolute Deviation.
	
	\[
	M_i = \frac{0.6745(x_i - median(x))}{MAD}
	\]
	
	where:
	
	\[
	MAD = median(|x_i - median(x)|)
	\]
	
	A value is flagged as suspicious when:
	
	\[
	|M_i| > 3.5
	\]
	
	This method is useful for Formula 1 data because lap times and speed values may be skewed by pit stops, special laps, or race interruptions.
	
	% ------------------------------------------------------------
	\subsection{Consensus Strategy}
	% ------------------------------------------------------------
	
	No single outlier detection method is perfect.
	For this reason, the script uses a simple consensus strategy.
	
	Each method can flag a value.
	The number of methods that flag the same value is counted.
	
	\begin{longtable}{p{3cm}p{5cm}p{6cm}}
		\toprule
		\textbf{Consensus Score} & \textbf{Meaning} & \textbf{Action} \\
		\midrule
		\endfirsthead
		
		\toprule
		\textbf{Consensus Score} & \textbf{Meaning} & \textbf{Action} \\
		\midrule
		\endhead
		
		0
		& Not flagged by any method.
		& Normal value. \\
		\midrule
		
		1
		& Flagged by one method.
		& Weak anomaly. Review only if relevant for the analysis. \\
		\midrule
		
		2
		& Flagged by both IQR and Modified Z-score.
		& Stronger suspicious value. Investigate before using the row in sensitive analysis. \\
		\bottomrule
	\end{longtable}
	
	The script does not remove consensus outliers automatically.
	It only marks them for later review.
	
	
	% ============================================================
	\section{Outlier Detection on the Lap Table}
	% ============================================================
	
	The main outlier detection activity is applied to the \texttt{lap} table.
	
	The selected numerical attributes are:
	
	\begin{table}[H]
		\centering
		\small
		\begin{tabular}{p{4cm}p{4cm}p{4cm}}
			\toprule
			\textbf{Timing Attributes} & \textbf{Timing Attributes} & \textbf{Speed Attributes} \\
			\midrule
			\texttt{lap\_time\_ms}     & \texttt{sector3\_time\_ms} & \texttt{speed\_i1} \\
			\texttt{sector1\_time\_ms} &                         & \texttt{speed\_i2} \\
			\texttt{sector2\_time\_ms} &                         & \texttt{speed\_fl} \\
			&                         & \texttt{speed\_st} \\
			\bottomrule
		\end{tabular}
	\end{table}
	
	Outliers are detected within each \texttt{session\_id}, but the analysis is not
	applied to all laps of the session.
	
	This is important because lap times and speed values are strongly affected by the
	racing context. A slow lap is not necessarily a data error. It may be caused by
	rain, Safety Car, Virtual Safety Car or other racing situations.
	
	For this reason, the script applies outlier detection only to laps that are suitable
	for performance comparison.
	
	% ------------------------------------------------------------
	\subsection{Clean Lap Selection}
	% ------------------------------------------------------------
	
	Before applying outlier detection, the script selects only clean candidate laps.
	
	A lap is considered a clean candidate for outlier detection when the following
	conditions are satisfied:
	
	\begin{table}[H]
		\centering
		\small
		\begin{tabular}{p{12cm}}
			\toprule
			\textbf{Clean Candidate Conditions for Outlier Detection} \\
			\midrule
			\texttt{lap\_time\_ms} is present \\[0.1cm]
			\texttt{compound} is present \\[0.1cm]
			\texttt{pit\_in\_time\_ms} is null \\[0.1cm]
			\texttt{pit\_out\_time\_ms} is null \\[0.1cm]
			\texttt{track\_status = 1} (AllClear) \\[0.1cm]
			The corresponding session is classified as dry \\
			\bottomrule
		\end{tabular}
	\end{table}
	
	The goal of this filter is not to clean the data.
	The goal is to avoid detecting obvious race-context effects as statistical outliers.
	
	For example, a lap under Safety Car would probably be detected as an outlier because
	it is much slower than normal laps. However, this does not mean that the value is
	wrong. It means that the lap belongs to a different racing context.
	
	% ------------------------------------------------------------
	\subsection{Dry Session Selection}
	% ------------------------------------------------------------
	
	Weather information is stored in the \texttt{weather} table as time-varying
	observations associated with each session.
	
	A precise approach would align each lap with the closest weather observation using
	\texttt{session\_id} and \texttt{time\_ms}. However, this would make the first
	version of the focused script more complex.
	
	For this project, the script uses a simpler and more conservative rule:
	
	\begin{center}
		\textit{A session is considered dry only if all weather observations for that
			session have \texttt{rainfall = false}.}
	\end{center}
	
	Therefore, if at least one weather observation in a session has
	\texttt{rainfall = true}, the whole session is excluded from lap outlier detection.
	
	This choice may exclude some usable dry laps from mixed sessions, but it avoids
	incorrectly comparing dry laps and wet laps in the same statistical group.
	
	% ------------------------------------------------------------
	\subsection{Track Status Filter}
	% ------------------------------------------------------------
	
	The \texttt{lap} table already contains a \texttt{track\_status} attribute.
	In this project, the outlier detection script keeps only laps with:
	
	\[
	\texttt{track\_status = 1}
	\]
	
	where status code \texttt{1} means \texttt{AllClear}.
	
	Laps with Safety Car, Virtual Safety Car, yellow flag, red flag, or other abnormal
	track conditions are excluded from the outlier detection sample.
	
	This is necessary because these laps are expected to have abnormal lap times, sector
	times, or speed values for domain reasons.
	
	% ------------------------------------------------------------
	\subsection{Session-Level Outlier Detection}
	% ------------------------------------------------------------
	
	After applying the clean lap, dry session, and track status filters, the script groups
	the remaining laps by \texttt{session\_id}.
	
	For each selected metric, the script applies the IQR method and the Modified Z-score
	inside each session group.
	
	The comparison is performed within the same session because lap times and speeds are
	strongly dependent on the circuit and session context. Comparing all laps across the
	entire dataset would be misleading. For example, Monaco and Monza naturally have very
	different lap times and speed profiles.
	
	The script applies the outlier methods only when a session contains a sufficient number
	of valid observations for the selected metric. If too few values are available, the
	check is skipped for that metric and session.
	
	The output includes:
	
	\begin{table}[H]
		\centering
		\small
		\begin{tabular}{p{12cm}}
			\toprule
			\textbf{Recorded Information for Each Outlier Candidate} \\
			\midrule
			Row identifier \\[0.1cm]
			Session identifier \\[0.1cm]
			Metric name \\[0.1cm]
			Observed value \\[0.1cm]
			IQR method flag \\[0.1cm]
			Modified Z-score flag \\[0.1cm]
			Consensus score \\[0.1cm]
			Interpretation label \\
			\bottomrule
		\end{tabular}
	\end{table}
	
	% ------------------------------------------------------------
	\subsection{Interpretation}
	% ------------------------------------------------------------
	
	A flagged lap value does not automatically mean that the value is wrong.
	
	Even after filtering, an outlier may still be caused by domain reasons such as traffic,
	driver mistakes, tyre degradation, unusual driving strategy, or source-specific
	recording behavior.
	
	For this reason, outlier rows are only flagged.
	The final decision is postponed to the Data Cleaning or analytical filtering phase.
	
	The focused script therefore produces evidence, not automatic corrections.

	% ============================================================
	\section{Expected Output Files}
	% ============================================================
	
	The focused script is expected to produce four main output files.
	
	\begin{longtable}{p{6cm}p{9cm}}
		\toprule
		\textbf{Output File} & \textbf{Purpose} \\
		\midrule
		\endfirsthead
		
		\toprule
		\textbf{Output File} & \textbf{Purpose} \\
		\midrule
		\endhead
		
		\texttt{focused\_missing\_summary.csv}
		& Summary of missing values by table and column. It reports the number and percentage of missing values and the main interpretation. \\
		\midrule
		
		\texttt{focused\_missing\_row\_flags.csv}
		& Row-level missing value flags. It reports the table, row identifier, column name, missing class, explanation, and severity. \\
		\midrule
		
		\texttt{lap\_outlier\_flags.csv}
		& Row-level outlier flags for lap-level metrics. It reports the metric, value, IQR flag, Modified Z-score flag, and consensus score. \\
		\midrule
		
		\texttt{outlier\_summary.csv}
		& Aggregated summary of detected outliers by table and metric. \\
		\bottomrule
	\end{longtable}
	

	
\end{document}